\section{A second domain, and where the law stops}
\label{sec:domain}

Every measurement so far is a Rubik's-cube sub-problem, so every claim above is
one state space seen through several sub-problems of itself. This section adds
the sliding-tile puzzle: the $3\times3$ board, which is the 8-puzzle, and the
$2\times4$. Exact distances come from breadth-first search over a Lehmer ranking
of the board. The implementation is checked against numbers it cannot influence:
it reproduces $n!/2$ reachable states on every board, the published diameters of
6, 21, 36 and 31 for $2\times2$, $2\times3$, $2\times4$ and $3\times3$, and the
known count of exactly two states at distance 31 on the 8-puzzle.

Abstractions are built as they are on the cube: track the blank and tiles
$1..j$, forget the rest. The measurement code is not duplicated. Both domains go
through one implementation of shell binning, pair sampling and tie counting,
because two scripts that merely resemble each other would let any difference
between them surface as a domain effect.

One trap is worth recording, because it would have flattered the classical
baseline rather than harming it. On the 8-puzzle, $j = 6$ leaves two tiles
untracked, and the parity invariant then fixes which way round those two go, so
the abstraction still separates every state and its table \emph{is} the exact
distance. As a baseline it scores $1.000$ at every shell and a GDRC of exactly
one. That is an oracle, not a strong heuristic. We exclude any rung whose
reachable count has not fallen below the full task's.

\paragraph{The profile results replicate.} On the 8-puzzle the learned heuristic
reaches GDRC $+0.965$, with adjacent-shell accuracy falling from $1.000$ to
$0.711$ across the measured shells. The nearest abstraction in strength, $j = 5$
at GDRC $+0.995$, falls from $1.000$ to $0.690$: slightly stronger overall and
decaying slightly \emph{more}. The random control stays flat at chance. Decay
with distance is not a fingerprint of bootstrapping in this state space either,
which is the claim of Section~\ref{sec:profile} surviving a change of domain.

Repeating that section's matched-strength control here, the abstraction decays
more than the learned heuristic in all 4 of the nearest-strength pairs available,
at a median GDRC gap of $0.030$ and a mean decay of $0.351$ against $0.296$. Four
pairs is too few to attach a significance test to, and the four learned models
are near-identical in decay, so we report the count and stop there. Note also
that decay is not comparable between the two domains: it is the first measured
shell's accuracy minus the last, and the tile is profiled over 20 shells against
the cube's 7, so the control is run within each domain and never pooled.

\paragraph{The law is worse here.} Split-sample mean absolute error is $0.0117$
over 252 cube observations and $0.0429$ over 239 sliding-tile observations,
against an estimator noise floor near $0.002$ in both. Read as a domain split
that says Equation~\ref{eq:law} does not generalise, and we first wrote it down
that way.

\paragraph{It is not a domain fact.} What the error tracks is the skewness of the
within-shell value distribution, Spearman $+0.74$, and the relationship is
monotone across all three tasks (Table~\ref{tab:domain}). Matched on skew, the
domain gap disappears below $|\mathrm{skew}| = 1$, where the tile fits slightly
\emph{better} than the cube in every bin. Reweighting the cube to the tile's skew
composition moves it only from $0.0221$ to $0.0237$, nowhere near the tile's
$0.0595$, which says the cube and the tile are not two populations with different
laws but one population indexed by skew.

\begin{table}[t]
\centering
\begin{tabular}{lrrrr}
\toprule
task & $n$ & median $|$skew$|$ & mean $|$error$|$ & bias \\
\midrule
cube, $k\!=\!6$ & 24 & 0.599 & 0.0221 & -0.0220 \\
sliding tile, $3\times3$ & 64 & 0.858 & 0.0508 & +0.0465 \\
sliding tile, $2\times4$ & 46 & 1.557 & 0.0718 & +0.0666 \\
\midrule
\multicolumn{5}{l}{\emph{Matched on skew, the domain gap closes below $|\mathrm{skew}| = 1$:}} \\
$|$skew$|$ range & cube $n$ & cube $|$err$|$ & tile $n$ & tile $|$err$|$ \\
0 to 0.4 & 5 & 0.0089 & 29 & 0.0065 \\
0.4 to 0.7 & 9 & 0.0199 & 12 & 0.0138 \\
0.7 to 1 & 2 & 0.0272 & 15 & 0.0243 \\
1 to $\infty$ & 8 & 0.0315 & 54 & 0.1080 \\
\midrule
\multicolumn{5}{l}{Reweighting the cube to the tile's skew mix gives 0.0237, against the tile's 0.0595.} \\
\bottomrule
\end{tabular}
\caption{The law's error tracks within-shell skewness rather than the state
space. Top: the three tasks in increasing order of median skew. Bottom: matched
on skew, the cube and the sliding tile agree below $|\mathrm{skew}| = 1$ and
diverge above it. Pattern databases only, so this reproduces without a trained
model.}
\label{tab:domain}
\end{table}

The class breakdown says the same thing from the other side
(Table~\ref{tab:domainclass}). Moving to the sliding tile costs learned
heuristics a factor of $1.9$ in mean absolute error and pattern databases a
factor of $3.3$. A learned value function is smooth and close to symmetric
within a shell; a coarse abstraction is not, and the sliding tile's abstractions
are the more skewed of the two. That is why the cube alone could not have
located this boundary: its abstractions never reach the skew where the law
begins to fail.

\begin{table}[t]
\centering
\begin{tabular}{lrrrr}
\toprule
 & \multicolumn{2}{c}{cube} & \multicolumn{2}{c}{sliding tile} \\
\cmidrule(lr){2-3}\cmidrule(l){4-5}
heuristic class & $n$ & mean $|$error$|$ & $n$ & mean $|$error$|$ \\
\midrule
learned & 138 & 0.0068 & 74 & 0.0126 \\
PDB & 80 & 0.0200 & 129 & 0.0667 \\
random & 34 & 0.0119 & 36 & 0.0203 \\
\bottomrule
\end{tabular}
\caption{Where the domain gap lives. Every class is worse on the sliding tile,
but pattern databases are worse by nearly twice the factor learned heuristics
are, which is the split the skew account predicts.}
\label{tab:domainclass}
\end{table}

\paragraph{Four explanations we tested and rejected.} Three are backwards from
the prediction, and we record them because a ruled-out explanation is the
cheapest thing to hand a reader who would otherwise reach for it first.
\emph{Saturation}, that a weak abstraction tops out at its own small diameter:
predicted the weakest rungs would misfit worst, measured Spearman $+1.000$ the
other way, the strongest rung being the worst fit at $0.122$ and the weakest the
best at $0.035$. \emph{Near-determinism}, that a within-shell spike is the least
normal shape available: error \emph{falls} with $d'$, Spearman $-0.53$, reaching
$0.0003$ above $d' = 5$. \emph{A strength confound}, which is this paper's own
warning about unmatched baselines turned on itself and so had to be checked:
reweighting the cube to the tile's $d'$ composition gives $0.0160$ against the
tile's $0.0595$, and within matched $d'$ bins the tile remains several times
worse. \emph{Ties}, which the equal-variance binormal identity assumes away and
which Section~\ref{sec:pooled} shows biases $\tau$-b: backwards again, since the
cube has a median tie rate of $0.325$ and the tile $0.000$.

Equal variance, the assumption we named explicitly when deriving
Equation~\ref{eq:law}, is not the culprit either. The tile's variance ratio is
$1.258$ against the cube's $1.428$, better behaved, and reweighting on it closes
almost none of the gap. The assumption that fails is normality, and specifically
its third moment, which we had not singled out.

\paragraph{What remains unexplained.} Above $|\mathrm{skew}| = 1$ the tile is
still $3.4$ times worse than the cube at matched skew. Skewness accounts for the
gap over most of the range and leaves a residual we cannot attribute. We report
the boundary we can demonstrate, not a mechanism we cannot.
